\documentclass{article}
\usepackage{hyperref}
\usepackage{datetime2}
\usepackage{textcomp}
\usepackage[left=0.2in, right=0.2in, top=0.2in, bottom=0.2in]{geometry}
\hypersetup{
    colorlinks=true,
    linkcolor=black,
    filecolor=magenta,
    urlcolor=black,
    pdftitle={Overleaf Example},
    pdfpagemode=FullScreen,
    }

\begin{document}

\begin{center}
    Marvin Wu(Wu Hung Mao)

hongmao.wu@mail.utoronto.ca $\mid$
\href{https://www.linkedin.com/in/hungmao-wu/}{Linkedin} $\mid$
\href{https://github.com/wuhungmao}{Github} $\mid$
(+1)437-970-3837 $\mid$
\href{https://wuhungmao.github.io/academic-work/}{Personal website}
\end{center}

\section*{Education}
\vspace{-8px}
Honours Bachelors of Science (H.Bsc), Computer Science specialist. \hfill Sep. 2021 - Present\\
University of Toronto Mississauga (expected graduation with PEY CO-OP - August 2026) \hfill GPA: 3.45/4.0
\vspace{-12px}

\section*{Technical Skills}
\begin{itemize}
    \item \textbf{Languages:} Python, C, C++, Java, SQL, JavaScript, HTML/CSS, PHP
    \vspace{-8px}
    \item \textbf{Frameworks \& Libraries:} React, Node.js, jQuery
    \vspace{-8px}
    \item \textbf{Tools \& Infrastructure:} Linux, Git/GitHub/GitLab, Docker, AWS, Jenkins, Jira
    \vspace{-8px}
    \item \textbf{Databases \& Profiling:} PostgreSQL, SQLite, Nsight Systems, LTTng, DCGM, Valgrind
    \vspace{-8px}
    \item \textbf{Other:} Agile/Scrum Methodologies, Spoken Languages (English, Mandarin)
\end{itemize}

\section*{Third/Fourth Year Courses}
\begin{itemize}
    \item \textbf{AI \& Robotics:} Introduction to Machine Learning, Introduction to Artificial Intelligence, Fundamentals of Robotics
    \vspace{-8px}
    \item \textbf{Systems \& Security:} Operating Systems, Parallel Programming, Principles of Computer Networks, Computer Security
    \vspace{-8px}
    \item \textbf{Software \& Data:} Introduction to Software Engineering, Programming on the Web, Introduction to Databases
    \vspace{-8px}
    \item \textbf{Theory \& Algorithms:} Algorithm Design and Analysis, Computational Complexity and Computability, Principles of Programming Languages
    \vspace{-8px}
    \item \textbf{Advanced Topics:} Computer Science Implementation Project, Topics in Computer Science, Computing Education
\end{itemize}

\section*{Work Experience}
\begin{itemize}
  \item[] \textbf{Physical Layer (vDU) Developer Intern} \hfill May 2024 - Sep 2025 \\
  \textit{Ericsson} \hfill \textit{Ottawa, Ontario, Canada}
  \vspace{-5px}
  \begin{itemize}
    \item[] \textbf{Project: Gphy2.0} (Feb 2025 - Present)
    \begin{itemize}
      \item Offloaded LDPC encoding onto the latest Nvidia GPU, previously running on an Intel accelerator.
      \item Created a new GPU branch for LDPC encoding in the existing codebase.
      \item Profiled latency of the GPU branch and compared Nsight System and Nsight Compute reports collected across servers with different GPUs, making adjustments to kernels to improve SM occupancy.
      \item Collected latency with the Intel accelerator branch calculated with the CPU cycle.
      \item \textbf{Environment:} Bazel\texttt{(rules\_cuda)}, Gerrit, Nsight System, Nsight Compute
    \end{itemize}
    \vspace{3px}
    \item[] \textbf{Project: Gemini} (May 2024 - Feb 2025)
    \begin{itemize}
      \item Digital twin for Radio, RF Channel, and UEs.
      \item Created a constellation diagram using the Ericsson Design System library and Vue3 framework.
      \item Wrote multi-threaded unit tests in CUDA to test a real-time system under a strict timing constraint ($\sim$ 35 \textmu s).
      \item Wrote an SQLite script that drew histograms based on data collected from LTTng traces for channel quality analysis.
      \item \textbf{Environment:} Docker, Wireshark, Gerrit, Gitlab, Jenkins, Jfrog, Sqlite3, AWS(EC2)
    \end{itemize}
  \end{itemize}
\end{itemize}

\section*{Academic Papers}
\begin{itemize}
\item[] Course Research, University of Toronto Mississauga \hfill Apr 2026
 \begin{itemize}
 \item[] \textbf{Wu, H.-M.}, Bi, F., Abdelhadi, Y., Lin, W., Khan, L. \href{https://drive.google.com/file/d/1KaVo7DM9onu-ap9cOaXpPX9bg1vr4znx/view?usp=sharing} {Ensemble Deep Learning Approaches for AI-Altered Video Detection}. Computer Science Implementation Project (CSC492).
    \vspace{3px}
    \item[] Riku, T., Ahmed, K., \textbf{Wu, H.-M.}, Usman, R. \href{https://drive.google.com/file/d/1YPqEK_rhctJtD1Wj-wzSBSfd9r5xmhBI/view?usp=sharing}{Can Agents Debug Better Than a Software Engineer?}. Course Research White Paper (CSC398).
  \end{itemize}
\end{itemize}

% ── AI & Machine Learning ──────────────────────────────────────────────────────
\section*{AI \& Machine Learning Projects}
\begin{itemize}

\item[] \textbf{AI-Video Detection Ensemble} $\mid$ \textit{Python, PyTorch, OpenCV, NumPy} \hfill \href{https://github.com/utmgdsc/AI-Video-Detection}{Link} $\mid$ Feb 2026
  \begin{itemize}
    \item Built an ensemble composed of 4 models (EfficientNet, MesoNet, XceptionNet, AASIST) to detect general \& deepfake AI-generated video.
    \item Fine-tuned a pre-trained EfficientNet-B1 model on datasets such as AIGVDBench and FaceForensics++, which use diffusion and face transfer methods to generate deepfake videos.
    \item Leveraged Transfer Learning to detect AI-generated media with high precision.
    \item Integrated the EfficientNet-B1 inference logic into the ensemble, utilizing MTCNN for face detection and cropping, and Albumentations for high-performance preprocessing.
    \item Designed an ensemble architectural diagram including preprocessing, inference, and postprocessing stages.
    \item Enabled parallel inference and took weighted scores to make the final REAL/FAKE prediction; communicated backend results with the frontend through HTTP requests.
    \item Co-authored the unpublished research manuscript {``Ensemble Deep Learning Approaches for AI-Altered Video Detection,''} detailing the multimodal architecture, performance metrics, and ensemble strategies.
    \item Served as both Project Manager and Senior Developer of a 5-member team; actively participated in two weekly meetings --- a Sunday sync facilitated by the mentor to review weekly progress, identify blockers, and plan upcoming tasks, and a Monday meeting to present results to the professor.
  \end{itemize}
  \vspace{3px}

  \item[] \textbf{JetBrains LCA Benchmark Contribution} $\mid$ \textit{Python, Docker, OpenAI API, Hugging Face API} \hfill \href{https://github.com/CSC392-CSC492-Building-AI-ML-systems/JetBrains-LCA-Winter2026}{Link} $\mid$ Feb 2026
  \begin{itemize}
    \item Contributed to the open-source JetBrains LCA benchmark, expanding the bug localization benchmark's capability from identifying the buggy file to automatically applying a fix.
    \item Used the OpenAI API to generate code patches and the Hugging Face API to programmatically download and manage large-scale repositories.
    \item Expanded evaluation metrics of the original benchmark by adding token usage count, iteration count, etc. for popular coding agents like Claude Code and Codex.
    \item Collaborated in a five-person team, aiming to support the long-term development of JetBrains' IDE-integrated agentic coder, Junie.
    \item Co-authored a technical white paper detailing the multi-agent architecture and evaluation metrics used to benchmark Claude Code and Codex.
    \item Delivered a group presentation to \textbf{40+ attendees} at the end of the project, communicating the multi-agent architecture, evaluation methodology, and benchmark findings.
  \end{itemize}
  \vspace{3px}

  \item[] \textbf{Machine Learning Painting Classifier} $\mid$ \textit{Python, Pandas, NumPy} \hfill \href{https://github.com/wuhungmao/Machine-Learning-Challenge}{Link} $\mid$ Feb 2026
  \begin{itemize}
    \item Used Pandas and NumPy to explore the dataset and analyze text features, identifying patterns to help the model distinguish between different paintings.
    \item Built a lightweight prediction script using only pure Python and NumPy because no external ML libraries are allowed, keeping the final model under 10MB for efficient deployment.
    \item Planned to test and tune 3 different types of models, selecting the best hyperparameters to ensure the model performs well on unseen test data rather than just memorizing the training set.
  \end{itemize}
  \vspace{3px}

  \item[] \textbf{Frozen Lake Solver} $\mid$ \textit{Python} \hfill \href{https://github.com/wuhungmao/frozen_lake}{Link} $\mid$ Feb 2025
  \begin{itemize}
    \item Implemented a reinforcement learning agent to solve the Frozen Lake environment from OpenAI Gym.
    \item Adopted the value iteration algorithm to find the optimal policy.
  \end{itemize}
  \vspace{3px}

  \item[] \textbf{Pacman AI Agent} $\mid$ \textit{Python} \hfill \href{https://github.com/wuhungmao/academic-work/tree/main/CSC384/Pacman\%20project/search}{Link} $\mid$ Sep 2023 - Oct 2023
  \begin{itemize}
    \item Implemented a Pacman game agent utilizing classical search algorithms (DFS, BFS, UCS, A*) with custom heuristics for the CSC384 AI course.
    \item Scored 80/100 on the assignment.
  \end{itemize}

\end{itemize}

% ── Systems & High-Performance Computing ──────────────────────────────────────
\section*{Systems \& High-Performance Computing Projects}
\begin{itemize}

  \item[] \textbf{Image Processing with GPU} $\mid$ \textit{C, CUDA} \hfill \href{https://github.com/wuhungmao/academic-work/tree/main/CSC367/Image\%20processing\%20with\%20GPU}{Link} $\mid$ Nov 2023 - Dec 2023
  \begin{itemize}
    \item Used GPU and parallel programming to process a 10MB x 10MB PGM image in under 10 milliseconds, achieving a 70x speedup compared to sequential processing. Scored 100/100 on the assignment.
    \item Broke down the entire project into smaller components, demonstrating strong time management skills that led to well-formatted code and detailed documentation.
    \item Employed C and CUDA to parallelize execution.
  \end{itemize}
  \vspace{3px}

  \item[] \textbf{Database Join with OpenMP and MPI} $\mid$ \textit{C, OpenMP, MPI} \hfill \href{https://github.com/wuhungmao/academic-work/tree/main/CSC367/Database\%20join\%20with\%20openMP\%20and\%20MPI}{Link} $\mid$ Oct 2023 - Nov 2023
  \begin{itemize}
    \item Performed database joins in both shared and distributed memory spaces, optimizing join speed. Achieved the merging of over ten thousand queries in under 200 milliseconds.
    \item Utilized C along with parallel programming concepts including OpenMP for multi-core and MPI for multi-node execution.
  \end{itemize}
  \vspace{3px}

  \item[] \textbf{File System Simulation} $\mid$ \textit{C, mmap, Pthreads} \hfill \href{https://github.com/wuhungmao/academic-work/tree/main/CSC369/File_Systems}{Link} $\mid$ Sep 2025 - Dec 2025
  \begin{itemize}
    \item Led a team to implement a Unix-like file system driver in C, utilizing memory mapping (mmap) to directly manipulate disk images.
    \item Managed disk layout structures including superblocks, inode tables, and bitmaps to track free inode and data blocks. Allocated resources as needed.
    \item Designed helper functions for path resolution and directory entry parsing to support nested directories and file lookups.
    \item Engineered logic for hard and symbolic links, managing reference counts and inode pointers to ensure data consistency.
    \item Designed a thread-safe synchronization scheme using fine-grained mutex locks to prevent race conditions during concurrent file operations.
  \end{itemize}
  \vspace{3px}

  \item[] \textbf{Ethernet Frame and ARP Request Simulation} $\mid$ \textit{C++} \hfill \href{https://github.com/wuhungmao/academic-work/tree/main/CSC358/pa1_starter_code}{Link} $\mid$ Feb 2024
  \begin{itemize}
    \item Implemented a \texttt{NetworkInterface} in C++ that resolves IP-to-MAC addresses via ARP requests and caches replies, simulating Layer 2 switch forwarding.
    \item Demonstrated understanding of Ethernet frame encapsulation, ARP protocol flow, and network switch behaviour.
  \end{itemize}

\end{itemize}

% ── Web & Full-Stack Development ──────────────────────────────────────────────
\section*{Web \& Full-Stack Development Projects}
\begin{itemize}

  \item[] \textbf{Personal Portfolio Website} $\mid$ \textit{React 18, Vite, React Router, CSS Modules, GitHub Actions} \hfill \href{https://wuhungmao.github.io/academic-work/}{Link} $\mid$ Dec 2023 - Present
  \begin{itemize}
    \item Designed and built a personal portfolio at \href{https://wuhungmao.github.io/academic-work/}{wuhungmao.github.io/academic-work} with a terminal/hacker aesthetic using React 18 and Vite; features 7 pages including a typewriter home, project gallery, course reviews, publications, and an interactive timeline.
    \item Migrated from an original plain HTML/CSS/jQuery site (built from scratch to self-teach web development) to a component-based React architecture with CSS Modules for per-component scoped styling and React Router v6 for client-side navigation.
    \item Configured automated CI/CD via GitHub Actions: every push to \texttt{main} triggers a Vite build and deploys \texttt{portfolio/dist/} to GitHub Pages with no manual steps.
    \item Utilized vibe coding with Claude Code to rapidly iterate on features, styling, and content.
  \end{itemize}
  \vspace{3px}

  \item[] \textbf{Decidophobia.com} $\mid$ \textit{Django, React, JavaScript, HTML/CSS, SQL} \hfill \href{https://github.com/cfstar188/Decidophobia.com}{Link} $\mid$ Jan 2024 - Apr 2024
  \begin{itemize}
    \item Acted as a Full Stack Developer to develop a responsive website that compares products from Amazon, Best Buy, eBay, and others, tailoring search results based on user preferences.
    \item Contributed to weekly sprint and non-sprint meetings, actively engaging in group discussions and proposing achievable and impactful features for implementation.
    \item Successfully retrieved product information from eBay and displayed it on the website.
    \item Monitored project progress with Jira and utilized the Django framework following the MVC model. Designed a questionnaire page using HTML, JavaScript, and CSS. Implemented back-end functionality, including product model creation and SQL queries for database storage.
    \item Converted a pure HTML document to React to create a more responsive website.
  \end{itemize}
  \vspace{3px}

  \item[] \textbf{Microservice API} $\mid$ \textit{Java, SQLite} \hfill \href{https://github.com/wuhungmao/academic-work/tree/main/CSC301}{Link} $\mid$ Jan 2024 - Apr 2024
  \begin{itemize}
    \item Implemented a bidirectional REST-based API, handling HTTP requests on both server and client sides. The API communicated with the database to update User and Product tables, providing HTTP responses upon completion.
    \item Created this project from the ground up and self-taught numerous new concepts.
    \item Collaborated with team members to break down the entire project into smaller, manageable tasks.
    \item Employed Java to interact with an SQLite database, leveraging JSON-formatted messages.
  \end{itemize}
  \vspace{3px}

  \item[] \textbf{Game Collection Website} $\mid$ \textit{PHP, PostgreSQL} \hfill Feb 2024
  \begin{itemize}
    \item Developed a website which incorporates different kinds of games (rock paper scissors, guessing numbers, frog puzzle).
    \item Developed multiple games utilizing PHP solely, incorporating concepts such as session management, HTTP requests, the Model-View-Controller (MVC) architecture, CSRF token, etc. Game information was stored in a PostgreSQL database management system.
  \end{itemize}
  \vspace{3px}

  \item[] \textbf{Wordle Game} $\mid$ \textit{Node.js, Express.js, jQuery} \hfill Mar 2024
  \begin{itemize}
    \item Designed a Wordle game which allows the player to guess a 5-letter word.
    \item Utilized the jQuery library to run JavaScript code on the front end and the Express.js framework to build a web application with Node.js and construct a RESTful API.
  \end{itemize}

\end{itemize}

% ── Other Course Projects ──────────────────────────────────────────────────────
\section*{Other Course Projects}
\begin{itemize}

  \item[] \textbf{Boggle Game} $\mid$ \textit{Java, JavaFX} \hfill \href{https://github.com/wuhungmao/academic-work/tree/main/CSC207}{Link} $\mid$ Oct 2022 - Dec 2022
  \begin{itemize}
    \item Led a team using Scrum methodology for Boggle game development, fostering active team communication. Successfully decomposed large group tasks, assigned them, and coordinated team efforts, resulting in a high-grade achievement.
    \item Created user personas and user stories, and generated detailed reports, showcasing excellent project management.
    \item Developed the game using Java and JavaFX, while utilizing GitHub for software development and project management.
  \end{itemize}
  \vspace{3px}

  \item[] \textbf{Simon Game} $\mid$ \textit{RISC-V Assembly} \hfill \href{https://github.com/wuhungmao/academic-work/tree/main/CSC258}{Link} $\mid$ Feb 2023 - Mar 2023
  \begin{itemize}
    \item Developed a Simon game using assembly code, designed to enhance players' short-term memory, resulting in a score of 90/100.
  \end{itemize}
  \vspace{3px}

  \item[] \textbf{File Compression} $\mid$ \textit{Python} \hfill \href{https://github.com/wuhungmao/academic-work/tree/main/CSC148/a2\%20huffman\%20algorithm\%20and\%20compression}{Link} $\mid$ Jan 2022 - Feb 2022
  \begin{itemize}
    \item Used the Huffman algorithm for both file compression and decompression.
  \end{itemize}

\end{itemize}

\end{document}
